\section{Discussion}

\subsection{Separating recognition support from the requested answer}
The central result is that an output boundary need not be an input boundary.
Keeping neighboring visual evidence improves recognition of a fixed request,
and the real-versus-padding intervention shows that this evidence supplies more
than additional frames. The interval formulation makes that distinction usable
with existing recognizers: context is admitted at encoding, while output scope
can be enforced at the state or transcript level.

The matched comparisons locate the value of projection more precisely.
Contextual direct decoding captures most of the benefit of admitting context.
Projection adds a smaller improvement at the principal operating point in all
four corpora, with more substantial effects in several short-target conditions.
The LRS2 availability analysis shows why an aggregate over short clips can
understate this opportunity: most two-second requests already cover the whole
clip. Together with the decoder screen, the results support treating readout as
a design decision conditioned on the recognizer and request, rather than using
one output rule for every operating point.

\subsection{Practical use and computational reuse}
The intended application is passage-level visual transcription: an editor or
viewer selects a time span and receives the words assigned to it, with a visible
allowance for surrounding video. This is useful when visual input is the
available speech signal and the full transcript has not already been computed.
The present evidence measures recognition and computational cost; a user study
would be needed to quantify editing-time or accessibility benefits.

Two workload regimes follow from the measurements. An isolated request can
benefit from processing a bounded passage; repeated requests can amortize a
full-clip transcript and temporal index. Full-clip recognition also remains
stronger for complete transcripts. Within bounded processing, direct readout
offers the lower measured median cost, while projection trades additional
decoding and alignment for its accuracy benefit. These operating distinctions
give applications a concrete basis for selecting an inference strategy.

\subsection{Scope of the evidence}
Word timing is automatic. The tested reference perturbations preserve the
favorable C--A and C--B point-difference signs, including matched decoders,
but fine allocation is more sensitive: a \(+160\) ms shift makes uniform-word
allocation better than aligned projection by 1.041 points on VLR and 0.245 on
LRS2. Independent review of boundary words would clarify this finer comparison.
The four corpus/model settings also differ in duration, training history, and
source diversity; Spanish has 12 recording sources and high absolute error.
The assessment material was available in earlier research. Results therefore
support the stated within-corpus comparisons, with prospective source-held
replication and evaluation outside the present recognizer family as valuable
extensions. The cost measurements characterize a stratified local-GPU sample.

\subsection{Toward request-dependent computation}
The next methodological question is whether observable request geometry and
predictive uncertainty can select context or readout more efficiently than a
fixed strategy. Such a policy should be developed on separate sources and
evaluated against fixed context budgets, contextual direct, projection, and
cached full-clip lookup under compatible input permissions. Its feature
extraction and routing costs must be included. The current duration and workload
results provide the empirical basis for this question; they are not an
evaluation of an adaptive policy.

\section{Conclusion}
Context-aware interval transcription separates the evidence used to recognize
speech from the span of text requested by the user. Prediction-aligned projection
implements this separation using a pretrained visual recognizer and its native
CTC scores. Four-corpus experiments show improved requested-interval recognition
over input cropping, while matched decoders distinguish contextual encoding
from the incremental effect of transcript projection. Context availability and
computational reuse explain the practical operating conditions. This formulation
extends pretrained VSR to selective passage requests and provides an empirical
basis for balancing context, readout, and work across queries.
